% =============================================================
\section{Diskussion und Grenzen der Compton-Formel}
% =============================================================

Die Herleitung und experimentelle Bestätigung der Compton-Formel
\[
\Delta \lambda = \frac{h}{m_e c}(1 - \cos\theta)
\]
war ein Meilenstein der Physikgeschichte. Dennoch gilt sie nur innerhalb bestimmter Grenzen
und beruht auf idealisierten Annahmen. Dieses Kapitel diskutiert ihren Gültigkeitsbereich,
ihre Grenzen und ihre Bedeutung im Kontext moderner Theorien.

% -------------------------------------------------------------
\subsection{Gültigkeitsbereich}
% -------------------------------------------------------------

Die Compton-Formel beschreibt die elastische Streuung eines Photons an einem Elektron unter der Annahme, dass
\begin{itemize}
	\item das Elektron vor dem Stoß ruht,
	\item der Stoß elastisch ist (\(E_\text{Photon} + E_\text{Elektron} = \text{konstant}\)),
	\item und nur \emph{ein Photon} mit \emph{einem Elektron} wechselwirkt.
\end{itemize}

\begin{PhysikBox}[Gültigkeit der Compton-Formel]
	Die Formel gilt exakt nur für freie Elektronen.
	In realen Materialien sind Elektronen jedoch an Atome gebunden. 
	Ihre Bindungsenergien führen zu kleinen Abweichungen vom idealen Verlauf,
	die experimentell als \emph{Compton-Profil} beobachtet werden.
\end{PhysikBox}

% -------------------------------------------------------------
\subsection{Grenzen bei hohen Photonenergien}
% -------------------------------------------------------------

Für sehr kurzwellige Photonen, insbesondere im Gamma-Bereich (\(E_\gamma \gtrsim 1~\text{MeV}\)),
wird der Compton-Effekt von weiteren quantenmechanischen Prozessen überlagert:

\begin{itemize}
	\item Der Rückstoß des Elektrons wird relativistisch.
	\item Bei Energien über \(1{,}022~\text{MeV}\) kann Paarbildung (\(\gamma \to e^+ + e^-\)) auftreten.
	\item Der Streuquerschnitt hängt nun stark vom Winkel ab und wird durch die \emph{Klein–Nishina-Formel} beschrieben.
\end{itemize}

\begin{DidaktikBox}[Relativistische Erweiterung]
	Bei sehr hochenergetischen Photonen muss die klassische Compton-Formel
	durch die \emph{Klein–Nishina-Gleichung} ersetzt werden.
	Diese beschreibt den differentiellen Streuquerschnitt relativistisch korrekt
	und stammt aus der Quantenelektrodynamik (QED).
\end{DidaktikBox}

% -------------------------------------------------------------
\subsection{Grenzen bei niedrigen Photonenergien}
% -------------------------------------------------------------

Für langwellige Photonen (\(h\nu \ll m_e c^2\)) ist die Compton-Verschiebung so klein,
dass sie experimentell nicht mehr nachweisbar ist. In diesem Grenzfall
geht der Compton-Effekt in die klassische \emph{Thomson-Streuung} über.

\begin{HinweisBox}[Übergang zur Thomson-Streuung]
	Im Niedrigenergiebereich verhalten sich Photonen wie klassische elektromagnetische Wellen.
	Die Elektronen schwingen im Feld des Lichts und strahlen erneut –
	die Wellenlänge bleibt dabei praktisch unverändert.
\end{HinweisBox}

% -------------------------------------------------------------
\subsection{Theoretische Bedeutung und Grenzen des Modells}
% -------------------------------------------------------------

Die Compton-Formel ist ein überzeugender Beweis für den Teilchencharakter des Lichts,
doch sie beschreibt Photon und Elektron noch als klassische Stoßpartner.
In der modernen Quantenfeldtheorie wird die Streuung als Wechselwirkung
zweier quantisierter Felder beschrieben – des elektromagnetischen Feldes
und des Elektronenfeldes.

\begin{DidaktikBox}[Grenze des Wellen–Teilchen-Bildes]
	Der Compton-Effekt lässt sich anschaulich mit Photonenimpulsen erklären,
	doch das Konzept „Photon als Teilchen“ ist nur eine Näherung.
	In der Quantenelektrodynamik entsteht der Effekt durch Austausch virtueller Photonen
	zwischen Elektron und elektromagnetischem Feld.
\end{DidaktikBox}

% -------------------------------------------------------------
\subsection{Historische Bedeutung}
% -------------------------------------------------------------

\begin{HistoryBox}[Comptons Vermächtnis]
	Arthur H. Compton erhielt 1927 den Nobelpreis für Physik,
	weil er experimentell zeigte, dass Licht Quantenimpuls trägt.
	Seine Arbeiten bestätigten die duale Natur des Lichts und leiteten
	den Übergang von der klassischen Wellentheorie zur Quantentheorie ein.
\end{HistoryBox}

% -------------------------------------------------------------
\subsection*{Zusammenfassung}
% -------------------------------------------------------------

\begin{itemize}
	\item Die Compton-Formel gilt exakt für freie Elektronen bei elastischer Streuung.
	\item Für sehr hohe Energien müssen relativistische Korrekturen (Klein–Nishina) berücksichtigt werden.
	\item Für niedrige Energien verschwindet die Verschiebung – es bleibt die Thomson-Streuung.
	\item Die Formel markiert den Übergang von der klassischen Optik zur Quantenphysik.
\end{itemize}

\begin{PhysikBox}[Bedeutung für die moderne Physik]
	Der Compton-Effekt verbindet Licht, Materie und Impulserhaltung in einer einzigen Beziehung.
	Er war der erste direkte experimentelle Beweis für den Impuls des Photons
	und bildet bis heute eine Brücke zwischen klassischer Elektrodynamik
	und Quantenelektrodynamik.
\end{PhysikBox}
